\chapter{Memory Optimization}
\label{ch:ch21-memory-optimization}

\section{When Memory Limits Performance}
\label{sec:ch21-bottleneck}

Memory accesses can dominate execution cost when latency or bandwidth limits progress.
The cost depends on locality, hardware, and overlap with other work; it must be measured for the workload rather than assumed for every program.
\term{Memory optimization} attacks that wait from the compiler side.
It removes loads and stores the program does not need, forwards values directly from stores to loads, replaces memory traffic with registers, and arranges the remaining accesses into friendly patterns.
The earlier chapters built the pieces: redundancy in Chapter~\ref{ch:ch16-value-numbering-redundancy}, motion in Chapter~\ref{ch:ch17-code-motion}, loops in Chapters~\ref{ch:ch18-loop-analysis} and~\ref{ch:ch19-loop-transformations}, and the alias facts that make memory reasoning sound in Chapter~\ref{ch:ch20-alias-analysis}.
This chapter assembles those pieces into optimizations whose currency is memory traffic rather than operation counts.
The central theme is that every memory optimization is a value optimization guarded by an alias query.
If the compiler knows two references cannot overlap, loads become reusable values and stores become removable.
If it cannot prove separation, the memory operation stays.

\section{Removing Redundant Loads and Stores}
\label{sec:ch21-redundant-accesses}

Three transformations carry most of the weight.
\term{Redundant load elimination} replaces a load with a previously loaded value when the address must read the same cell.
\term{Store-to-load forwarding} replaces a load with the value of a dominating store to the same address.
\term{Dead store elimination} removes a store that no later load can observe.
Listing~\ref{lst:ch21-diamond} shows a diamond where the join block reloads a value available on both paths.

\begin{irlisting}[caption={A diamond where the join repeats a load from before the branch.},label={lst:ch21-diamond}]
entry: x = load p
       store q = 1
       br c, left, right
left:  y = x + 1
       goto join
right: z = x + 2
       goto join
join:  w = load p
       ret w
\end{irlisting}

The load in the join block is redundant when two conditions hold on every path.
First, the address \code{p} still holds the value loaded in the entry, which requires that no store on any path write a location that may overlap \code{p}.
The store to \code{q} is harmless exactly when the alias oracle answers no-alias for \code{p} and \code{q}.
Second, the load must be available on every path, meaning each path carries the value forward without killing it.
Here both paths carry \code{x} untouched, so the join load reduces to the copy \code{w = x}.
The analysis behind this decision is \term{available-load} dataflow, a forward analysis in the same family as the availability used for expressions.
A load is available at a point when every path computed that address with no intervening may-store to it.
Algorithm~\ref{alg:ch21-load-elim} computes availability and rewrites the redundant loads.

\begin{algorithm}
  \caption{Redundant load elimination with alias kills}\label{alg:ch21-load-elim}
  \begin{pseudocode}
    \Require reachable SSA CFG, ordinary nonvolatile non-atomic loads, sound effect and alias information
    \Ensure replace loads only with a dominating, unchanged value of matching address, type, and width
    \State $Avail(B) \gets$ empty set for every block $B$
    \State $changed \gets true$
    \While{$changed$}
      \State $changed \gets false$
      \For{each block $B$}
        \State $In \gets \bigcap_{P \in \Pred(B)} Avail(P)$; use the empty set at entry
        \State scan $B$ in order; kill clobbered facts and add each original load's typed address-value fact
        \State $New \gets$ facts remaining at the end of the scan
        \If{$New \neq Avail(B)$}
          \State $Avail(B) \gets New$; $changed \gets true$
        \EndIf
      \EndFor
    \EndWhile
    \State replay each block from its converged input, recording facts before each instruction
    \For{each load $l$ of address $a$ in block $B$}
      \If{a matching fact before $l$ names a value dominating $l$}
        \State replace $l$ with a copy of the recorded value
      \EndIf
    \EndFor
  \end{pseudocode}
\end{algorithm}

The algorithm assumes loads are keyed by the value number of their address, so equal addresses with equal numbers are comparable.
Facts also include the original load's SSA result, type, and access width; keep all such facts during solving rather than choosing a changing representative.
Starting from empty computes a conservative solution that can miss loop-carried availability.
It is a simple baseline, not the greatest must solution obtained by universal initialization over a fixed fact universe.
Replacement occurs only after convergence, uses facts before the candidate load, and updates uses so later replacements cannot refer to a deleted definition.
It assumes the alias oracle is sound: every store that might overlap a tracked address kills it.
Kills are deliberately coarse for calls, which remove every tracked load unless interprocedural summaries from Chapter~\ref{ch:ch22-interprocedural-analysis} say otherwise.
Dead store elimination runs the mirror analysis backward.
A store is dead only when its value cannot be observed before a complete overwrite or the end of the object's lifetime, and removing the access preserves required effects and faults.
Function return is not the end of observation for globals or escaped objects: their final values can be read by a caller.
The backward pass must seed those externally observable locations as live at exits and include possible reads by calls.
A may-write is not a guaranteed overwrite; killing backward memory liveness requires a must-write covering the relevant bytes.

\section{Worked Example: Availability Across the Diamond}
\label{sec:ch21-worked-availability}

We run the analysis of Algorithm~\ref{alg:ch21-load-elim} on Listing~\ref{lst:ch21-diamond} by hand, assuming the oracle reports no-alias for \code{p} and \code{q}.
Denote the load of \code{p} by $L$.
The entry block generates $L$ with value \code{x}: its input set is empty, it computes the load, and the store to \code{q} kills nothing because of the no-alias answer.
So $Avail(entry) = \{L\}$.
The left block generates nothing and kills nothing, passing its input through unchanged.
The right block behaves identically.
In the first iteration, the left and right blocks read the entry result, giving $Avail(left) = \{L\}$ and $Avail(right) = \{L\}$.
The join block meets its predecessors by intersection: $\{L\} \cap \{L\} = \{L\}$.
Its load of \code{p} therefore finds $L$ in its input set; after convergence it rewrites to \code{w = x}.
A second iteration changes nothing, so the fixed point is reached; the original join load also contributes its own result fact during solving.
Table~\ref{tab:ch21-avail} records the sets per block.

\begin{table}
  \centering
  \caption{Available-load sets for Listing~\ref{lst:ch21-diamond} with a no-alias answer for the store.}
  \label{tab:ch21-avail}
  \begin{tabular}{lll}
    \toprule
    Block & Input set & Output set \\
    \midrule
    entry & empty & $L$ \\
    left & $L$ & $L$ \\
    right & $L$ & $L$ \\
    join & $L$ & $L$, plus original join-load fact \\
    \bottomrule
  \end{tabular}
\end{table}

Now change one fact and watch the conclusion flip.
If the oracle instead answers may-alias for \code{p} and \code{q}, the store in the entry kills $L$ immediately.
The branch-arm outputs and join input become empty, so the join load stays.
That retained load generates a new fact at the join exit.
Figure~\ref{fig:ch21-diamond} draws the diamond with the store annotated by the oracle answer that decides the outcome.
One alias answer is the entire difference between removing a load and keeping it, which is why memory optimization budgets so much effort for pointer precision.

\begin{figure}
  \centering
  \begin{tikzpicture}[node distance=14mm and 20mm]
    \node[basic block] (ch21e) {E: x = load p; store q = 1};
    \node[basic block,below left=of ch21e] (ch21l) {L: y = x + 1};
    \node[basic block,below right=of ch21e] (ch21r) {R: z = x + 2};
    \node[basic block,below right=of ch21l] (ch21j) {J: w = load p};
    \draw[cfg edge] (ch21e) -- (ch21l);
    \draw[cfg edge] (ch21e) -- (ch21r);
    \draw[cfg edge] (ch21l) -- (ch21j);
    \draw[cfg edge] (ch21r) -- (ch21j);
  \end{tikzpicture}
  \caption{The join load is redundant only when the store cannot reach the loaded address, which is exactly what the alias oracle decides.}
  \label{fig:ch21-diamond}
\end{figure}

\section{Registers Instead of Memory}
\label{sec:ch21-registers}

Eliminating one load at a time wins battles, but the larger victory is promoting memory locations into registers wholesale.
\term{Scalar replacement} takes an array element or structure field accessed repeatedly in a loop and holds it in a temporary across iterations.
A loop that accumulates \code{a[0]} on every pass loads and stores the same cell each time.
After replacement, the loop works on a register temporary and writes the cell back once after the exit.
\term{Loop-invariant loads} are the motion counterpart: a load whose address never changes in the loop hoists to the preheader like any other invariant.
Promotion must account for every read and write that may observe the location: accesses replaced by the temporary are allowed, but unrewritten aliases and calls need proof or synchronization with memory.
Load hoisting additionally requires a safe execution point; an invariant address alone does not justify a faulting load on a zero-trip loop.
Delayed write-back must cover every relevant exit, including exceptional exits, and must not invent an access when the original loop executes zero times.
\term{Alignment and padding} handle the layout side.
A value split across a cache-line boundary costs extra transfers, and a structure with poorly ordered members wastes space on padding.
Reordering fields is legal only when layout is under compiler or runtime control and every dependent access is updated.
Published ABIs, C structure layout, address arithmetic, and serialization can make member offsets observable.
Table~\ref{tab:ch21-menu} collects the transformations with the fact each one needs.

\begin{table}
  \centering
  \caption{Memory optimizations with the fact that makes each one legal.}
  \label{tab:ch21-menu}
  \begin{tabular}{lll}
    \toprule
    Optimization & What it removes & Legality fact \\
    \midrule
    Load elimination & Repeated loads & Address available, no may-store \\
    Store forwarding & Load after store & Must-alias with nothing between \\
    Dead stores & Unread stores & Location not live on any path \\
    Scalar replacement & Loop traffic & All observers accounted for \\
    Invariant loads & Repeated hoists & Address invariant, safe to speculate \\
    Layout tuning & Wasted lines & Unobservable layout, updated uses \\
    \bottomrule
  \end{tabular}
\end{table}

Two language features restrict all of the above.
A \code{volatile} access in C has side effects the abstract machine defines, so the compiler must preserve its count and order relative to other volatile accesses and cannot remove or invent them~\cite{wg14-2011-c11}.
Atomic operations and synchronization likewise constrain motion across their boundaries, since reordering them would change what other threads can observe.
Treat volatile locations, atomics, and synchronization points as opaque barriers in every memory analysis.
Optimizers describe such locations as \enquote{hands off} memory, meaning no analysis may claim facts across them.
Calls are barriers too unless summaries describe them, which is the subject of Chapter~\ref{ch:ch22-interprocedural-analysis}.

\section{Implementation Notes}
\label{sec:ch21-implementation}

For $N$ blocks and $L$ original load facts, the tables need $O(NL)$ bits with a dense representation; this is a space bound, not the total running time.
Time includes repeated predecessor intersections, block scans, and alias queries until convergence.
Numbering addresses with value numbers keeps $L$ small by merging equal addresses before the analysis starts.
Dead-store elimination needs its own location universe, exit boundary, and byte-coverage checks before a comparable cost can be derived.
Scalar replacement is linear in the number of candidate references once alias answers exist, because each candidate needs one promotion check and one rewrite.
The real cost hides in the alias queries: each may-store asks the oracle about each live load, so caching answers per address pair matters.
Production compilers layer these passes: alias analysis first, then load elimination with forwarding, then dead stores, then scalar replacement inside loops, then cleanup numbering to merge the exposed values.
Memory dependence testing for loops adds a pairwise check over references scaled by dependence distance, which stays cheap for the inner loops that matter.
Cache behavior guides profitability rather than legality.
Blocking from Chapter~\ref{ch:ch19-loop-transformations} shapes the working set, prefetching hides latency for predictable streams, and both rely on the trip counts of Chapter~\ref{ch:ch18-loop-analysis}.
Measure with miss counters rather than static estimates, because a transformation that removes loads but doubles misses is a loss.

\begin{ccwarning}[Volatile means hands off]
Never eliminate, forward, or move a volatile access to save traffic.
The language defines its observable behavior, and the optimizer must preserve every such access exactly.
Fence volatile regions out of every memory analysis before tuning anything else.
\end{ccwarning}

\section{Exercises}
\label{sec:ch21-exercises}

\begin{enumerate}
  \item \label{ex:ch21-forward-vs-elim}
  Given the sequence \code{store p = 7}, \code{x = load p}, \code{store q = 2}, and \code{y = load p} with a must-alias answer for the first store and the first load and a no-alias answer for \code{p} and \code{q}, decide for each load whether forwarding or elimination applies.
  Rewrite the sequence, showing which stored value reaches each load.
  Then change only the second oracle answer to may-alias and redo the decision for the second load.
  \item \label{ex:ch21-dead-store}
  Draw a CFG where block $A$ stores to \code{m}, block $B$ stores to \code{m} again, and both reach an exit with no intervening load of \code{m} on one path but a load on the other.
  Compute live memory locations backward by hand and decide which store, if any, is dead.
  Explain why overwriting alone does not make the first store dead on the path that loads.
  \item \label{ex:ch21-scalar-replace}
  Write a loop that reads and writes \code{a[k]} each iteration for a loop-invariant index \code{k} and also reads \code{a[i]} for the varying index \code{i}.
  State the alias query that decides whether \code{a[k]} promotes to a register, perform the replacement under a no-alias answer, and show the write-back after the exit.
  Explain what breaks if a later edit adds a call inside the loop with unknown effects.
\end{enumerate}

\section{Further reading}
\label{sec:ch21-further-reading}

Morgan's \emph{Building an Optimizing Compiler} treats memory optimizations alongside the dataflow analyses that enable them~\cite{morgan-1998-optimizing}.
Read it for how load and store motion fits the book's wider optimization pipeline.
Appel's \emph{Modern Compiler Implementation in ML} shows the aliasing and liveness facts these passes consume~\cite{appel-1998-modern}.
Use it to place memory optimization correctly among its neighboring passes.
The C11 committee draft specifies volatile semantics, sequencing, and the observable behavior that bounds every transformation here~\cite{wg14-2011-c11}.
Check those clauses before removing or moving any access whose effects the program can observe.
